% =====================================================================
%  Chapter 31 opener -- "Chapter at a Glance", three reading paths + the
%  ladder staircase, THE running example (the 5-vertex WEIGHTED graph G on
%  {1,2,3,4,5} with edges {1-2:1, 1-3:3, 2-3:4, 3-4:2, 2-4:5, 4-5:6,
%  3-5:7} whose unique MST is {1-2, 3-4, 1-3, 4-5} of weight 12), the
%  cut-property / Kruskal / Prim / KRT / classification taxonomy, the two
%  traps (greedy needs the cut property not "locally cheapest", and the
%  Prim-is-Dijkstra confusion), the who-owns-what comparison against Ch.17
%  (DSU machinery -- the black box Kruskal calls), Ch.16 (LCA -- the KRT's
%  query engine), Ch.29 (Trees Born from Graphs -- two chapters back), and
%  Ch.30 (Counting Labelled Trees -- the PREVIOUS, enumeration twin), a
%  constraint -> technique table, and notation.
%  Exports fig:c17-tree, fig:c17-taxonomy.
%  \term the NEW objects "spanning tree", "minimum spanning tree", "cut
%  property", "Kruskal", "Prim", "Kruskal-reconstruction tree", "bottleneck
%  path" only -- DSU/LCA/depth are indexed by earlier chapters.
% =====================================================================

\begin{chaptertheorybox}{Chapter at a Glance}
Chapter~30 (\chref{c32}) ended the book's enumerative arc: \emph{how many}
labelled trees match a description.  This chapter is its optimisation twin.
The input is a single \emph{weighted} graph, and the question is no longer a
count but a \textbf{choice}: pick the \(n-1\) edges that connect everything at
least total cost.  That object -- the \term{minimum spanning tree} (MST) -- is
one of the two or three workhorses of all of competitive graph algorithms, and
the entire chapter is a meditation on one lemma that makes the greedy correct.

The method rests on a single idea, the \term{cut property}
(\(\secref{c17-kruskal}\)): for any partition of the vertices into two sides,
the cheapest edge crossing the partition belongs to some MST.  Every greedy
rule in this chapter is the cut property wearing a different hat.  \term{Kruskal}
(\(\secref{c17-kruskal}\)) sorts the edges and unions their endpoints with a DSU
(\chref{c28}), taking an edge exactly when it merges two components -- \(\bigO{m
\log m}\).  \term{Prim} (\(\secref{c17-prim}\)) grows one component with a
frontier priority queue -- the dense-graph choice.  Recording each Kruskal union
as a fresh node weighted by the causing edge yields the
\term{Kruskal-reconstruction tree} (\(\secref{c17-krt}\)), which converts
bottleneck path questions into LCA weights (\chref{c11}).  Finally
\(\secref{c17-classify}\) answers the sharper question -- \emph{which} edges can
appear in an MST -- via the cycle and intersection properties, splitting every
edge into forced, optional, or forbidden.

\textbf{You will learn to}
\begin{itemize}[nosep]
  \item state and apply the cut property, and prove Kruskal and Prim correct
        from it in two lines each (\(\secref{c17-kruskal}\)); \hfill\lv{3}
  \item run Kruskal with a DSU in \(\bigO{m\log m}\) and Prim with a heap in
        \(\bigO{m\log n}\), and choose between them by edge density
        (\(\secref{c17-prim}\)); \hfill\lv{4}
  \item build the Kruskal-reconstruction tree and answer \term{bottleneck}
        (minimax-edge) path queries as an LCA weight (\(\secref{c17-krt}\));
        \hfill\lv{4}
  \item classify edges as forced / optional / forbidden with the cycle property,
        and count equal-weight MST choices per Kruskal group
        (\(\secref{c17-classify}\)); \hfill\lv{5}
  \item compute a second-best MST by a single non-tree-edge swap, and recognise
        when a "spanning tree" question is really a Ch.29 robustness or a Ch.30
        counting problem (\(\secref{c17-boundary}\)). \hfill\lv{5}
\end{itemize}

\textbf{Prerequisites.} Chapter~17 (\chref{c28}) -- the DSU, used as a black box
by Kruskal (\texttt{find}/\texttt{union}); this chapter never re-derives it.
Chapter~16 (\chref{c11}) -- lowest common ancestor, the query engine the
Kruskal-reconstruction tree rides on.  Chapter~13 (\chref{c10}) -- the bottom-up
tree fold, reused by the Summit's divide-and-conquer over value bits.  Chapter~30
(\chref{c32}) -- the \textbf{previous chapter}; it ends at "here is how many
trees match a description"; this one opens at "here is the single cheapest one."
\end{chaptertheorybox}

% ---------------- Reading paths + ladder position ----------------
\begin{chaptertheorybox}{How to Read This Chapter (three paths)}
Level badges \(\lv3\)--\lv{5}: \(\lv3\) is stating the cut property and running
Kruskal by hand on the running graph; \(\lv4\) is coding Kruskal with a DSU,
Prim with a heap, and building the reconstruction tree for a bottleneck query;
\(\lv5\) is the MST edge classification (CF 1108F) and the Summit (CF 888G) -- a
complete graph whose edge weight is an xor, solved by a trie divide-and-conquer
that is the cut property over value bits.

\textbf{Greedy path (\(\lv3\)--\lv{4}, ``just build the MST'').}
\secref{c17-kruskal} in full -- cut property, Kruskal, solved CF 1095F -- then
\secref{c17-prim} for the dense-graph frontier view.  Stop here if your problem
only wants the MST weight or the tree itself.

\textbf{Query path (\(\lv4\)--\lv{5}, ``ask about paths in the MST'').}
\secref{c17-krt}: record the unions, and a bottleneck (max-edge-on-path) query
becomes an LCA weight (\chref{c11}).  This is the branch for problems that build
the MST \emph{and} then answer "cheapest threshold to connect \(u,v\)."

\textbf{Structure path (\(\lv5\), ``which edges, how many, how unique'').}
\secref{c17-classify} -- forced / optional / forbidden and the equal-weight group
count, capped by solved CF 1108F.  Read \(\secref{c17-boundary}\) after any
branch to know where the chapter's method stops and Ch.29 or Ch.30 resumes.
\end{chaptertheorybox}

\begin{figure}[H]
\centering
\begin{tikzpicture}[diagram, x=8.5mm, y=6.5mm]
  \foreach \st/\lab in {1/lv1,2/lv2,3/lv3,4/lv4,5/lv5,6/lv6} {
    \fill[\lab!22, draw=\lab, line width=0.8pt] (\st-1,0) rectangle (\st,\st*0.52);
    \node[font=\scriptsize\bfseries, text=\lab] at (\st-0.5,\st*0.52-0.42) {L\st};}
  \node[annot, anchor=south west, font=\scriptsize]
    at (0.05,0.62) {prereq: DSU (Ch.17), LCA (Ch.16), tree fold (Ch.13)};
  \node[annot, anchor=south west, font=\scriptsize]
    at (2.05,1.66) {31.1 cut property + Kruskal -- solved CF 1095F};
  \node[annot, anchor=south west, font=\scriptsize]
    at (3.05,2.18) {31.2 Prim frontier | 31.3 Kruskal-reconstruction tree};
  \node[annot, anchor=south west, font=\scriptsize, text=lv5]
    at (4.05,2.70) {31.4 edge classification -- solved CF 1108F};
  \node[annot, anchor=south west, font=\scriptsize, text=lv5]
    at (5.05,3.22) {31.5 boundary; Summit CF 888G (xor-MST, trie D\&C)};
\end{tikzpicture}
\caption{Chapter~31 on the six-step ladder: hand-running Kruskal on the running
graph is \(\lv3\); coded Kruskal/Prim and the reconstruction-tree bottleneck
query are \(\lv4\); the edge classification (CF 1108F) is \(\lv5\); the Summit
(CF 888G) sits at the top of step~5 for this part of the book.}
\end{figure}

% ---------------- The running example ----------------
\subsection*{The Running Example: A Weighted Graph on Five Vertices}

Every section up to the Summit runs on the same 5-vertex weighted graph \(G\),
fixed once here and authoritative (if a later trace disagrees, this figure
wins).  Vertices are \([5]=\{1,2,3,4,5\}\); the seven edges with their weights
are
\[
  1\!-\!2:1,\quad 1\!-\!3:3,\quad 2\!-\!3:4,\quad 3\!-\!4:2,\quad
  2\!-\!4:5,\quad 4\!-\!5:6,\quad 3\!-\!5:7 .
\]
Because all seven weights are distinct, \(G\) has a \emph{unique} MST, and
Kruskal, Prim, and the reconstruction tree must all recover the same four edges:
\(\{1\!-\!2,\ 3\!-\!4,\ 1\!-\!3,\ 4\!-\!5\}\) of total weight \(1+2+3+6=12\).
That single number is the checkpoint for every trace in the chapter.

\begin{figure}[H]
\centering
\begin{tikzpicture}[diagram, scale=0.98, transform shape,
   tnode/.style={circle, draw=primaryblue!70, fill=primaryblue!8, line width=0.8pt,
                 minimum size=8mm, font=\small\bfseries, inner sep=1pt},
   mst/.style={line width=1.7pt, draw=successgreen!80!black},
   rest/.style={line width=0.9pt, draw=darkgray!70}]
  \node[tnode] (v1) at (0,0) {1};
  \node[tnode] (v2) at (3,1.4) {2};
  \node[tnode] (v3) at (3,-1.6) {3};
  \node[tnode] (v4) at (6,-0.4) {4};
  \node[tnode] (v5) at (8.4,1.2) {5};
  \draw[mst] (v1)--node[fill=white,inner sep=1.5pt,font=\scriptsize] {1} (v2);
  \draw[mst] (v1)--node[fill=white,inner sep=1.5pt,font=\scriptsize] {3} (v3);
  \draw[mst] (v3)--node[fill=white,inner sep=1.5pt,font=\scriptsize] {2} (v4);
  \draw[mst] (v4)--node[fill=white,inner sep=1.5pt,font=\scriptsize] {6} (v5);
  \draw[rest] (v2)--node[fill=white,inner sep=1pt,font=\scriptsize] {4} (v3);
  \draw[rest] (v2)--node[fill=white,inner sep=1pt,font=\scriptsize] {5} (v4);
  \draw[rest] (v3)--node[fill=white,inner sep=1pt,font=\scriptsize] {7} (v5);
  \node[annot, font=\scriptsize, text width=8cm, align=center] at (4,-3)
    {green: the four MST edges (weight 12); gray: rejected by the cycle property};
\end{tikzpicture}
\caption{The running graph \(G\).  The thick green edges are its unique MST
\(\{1\!-\!2,3\!-\!4,1\!-\!3,4\!-\!5\}\), total weight \(12\).  The three gray edges
each close a cycle whose heaviest edge is itself, so the cycle property rejects
them.  This object is fixed for the whole chapter.}
\label{fig:c17-tree}
\end{figure}

Two traps return throughout.  The \textbf{"locally cheapest is globally safe"
trap}: a greedy that simply takes an edge when it looks cheap, without checking
it crosses a cut consistent with the edges already chosen, is wrong.  The cut
property is the \emph{only} permission a spanning-tree greedy needs; Kruskal and
Prim are two mechanical ways to honour it.  The \textbf{Prim-is-Dijkstra trap}
is subtler: both maintain a frontier priority queue, but Prim relaxes by the
edge's own weight to grow a tree, whereas Dijkstra relaxes by accumulated
distance from a source.  Feeding Dijkstra's rule to a spanning-tree problem
produces a shortest-path tree, not an MST.  Every confused implementation
conflates these two relaxations somewhere.

\subsection*{Which Job? Reading the Statement}

\begin{center}
\renewcommand{\arraystretch}{1.25}
\setlength{\tabcolsep}{3pt}
\begin{tabular}{@{}p{0.30\linewidth} p{0.20\linewidth} p{0.28\linewidth} >{\raggedright\arraybackslash}p{0.13\linewidth}@{}}
\toprule
\textbf{The statement says\dots} & \textbf{Tool} & \textbf{Why it works} & \textbf{Section}\\
\midrule
"cheapest set of edges to connect all"; sparse \(m\ll n^2\) & Kruskal + DSU & sort once, union per edge & \(\secref{c17-kruskal}\)\\
"connect all vertices"; dense graph / matrix given & Prim + heap & frontier relaxation, \(\bigO{n^2}\) & \(\secref{c17-prim}\)\\
"minimise the heaviest edge on the path between \(u,v\)" & reconstruction tree & path-max = LCA weight & \(\secref{c17-krt}\)\\
"is this edge in an MST?"; "how many MSTs"; "make MST unique" & classification & cycle / intersection property & \(\secref{c17-classify}\)\\
"second cheapest spanning tree"; "swap one edge" & replacement edge & non-tree edge closes a cycle & \(\secref{c17-boundary}\)\\
\bottomrule
\end{tabular}
\end{center}

\begin{insightbox}[The whole chapter in one sentence]
A spanning-tree greedy is correct exactly when each edge it takes is the
cheapest crossing of some cut the already-taken edges respect -- the cut
property -- and Kruskal, Prim, the reconstruction tree, and the forced/optional/
forbidden classification are four readings of the same lemma, so the entire craft
of this chapter is recognising which reading a problem is asking for and then
honouring the cut.
\end{insightbox}

% ---------------- Taxonomy figure ----------------
\begin{figure}[H]
\centering
\begin{tikzpicture}[diagram, scale=0.9, transform shape,
   mnode/.style={circle, draw=primaryblue!70, fill=primaryblue!8, line width=0.7pt,
                 minimum size=13mm, font=\scriptsize\bfseries, inner sep=1pt, align=center}]
  \node[mnode] (root) at (0,0) {MST via\\ cut property};
  \node[mnode] (kru)  at (-4,-1.7) {Kruskal\\ (sort + DSU)};
  \node[mnode] (pri)   at (0,-1.7) {Prim\\ (frontier)};
  \node[mnode] (krt)  at (4,-1.7) {recon-\\ struction};
  \node[mnode] (cla)   at (-2.4,-3.4) {edge\\ classify};
  \node[mnode] (sec)   at (2.4,-3.4) {second\\ best};
  \draw (root)--(kru)--(cla);
  \draw (root)--(pri);
  \draw (root)--(krt);
  \draw (kru)--(sec);
\end{tikzpicture}
\caption{The five readings of the cut property. A greedy \emph{build} is either
Kruskal (\(\secref{c17-kruskal}\)) or Prim (\(\secref{c17-prim}\)).  Recording
the build gives the reconstruction tree (\(\secref{c17-krt}\)); interrogating it
gives the classification (\(\secref{c17-classify}\)) and the second-best swap
(\(\secref{c17-boundary}\)).}
\label{fig:c17-taxonomy}
\end{figure}

\subsection*{Who Owns What}

\begin{center}
\renewcommand{\arraystretch}{1.2}
\setlength{\tabcolsep}{3pt}
\begin{tabular}{@{}p{0.28\linewidth} p{0.40\linewidth} p{0.24\linewidth}@{}}
\toprule
\textbf{Object} & \textbf{Owned here?} & \textbf{Elsewhere}\\
\midrule
Cut property, Kruskal, Prim & \textbf{Yes}, \(\secref{c17-kruskal}\)/\(\secref{c17-prim}\) & --\\
Kruskal-reconstruction tree & \textbf{Yes}, \(\secref{c17-krt}\) & --\\
MST edge classification, uniqueness count & \textbf{Yes}, \(\secref{c17-classify}\) & --\\
Second-best / replacement MST & \textbf{Yes}, \(\secref{c17-boundary}\) & --\\
DSU \texttt{find}/\texttt{union} & no -- used as a black box & Ch.17 (\chref{c28})\\
LCA (for KRT queries) & no -- used as a black box & Ch.16 (\chref{c11})\\
Bottom-up tree fold & no -- reused by the Summit & Ch.13 (\chref{c10})\\
Counting spanning trees of \(G\) & no -- a determinant, not a greedy & Ch.30 (\chref{c32})\\
Robustness / bridge questions on unweighted graphs & no & Ch.29 (\chref{c31})\\
\bottomrule
\end{tabular}
\end{center}

\subsection*{Recognition Guide}

\begin{recognitionbox}[Turning a graph into an MST question]
\reckey{``minimum total cost to connect everything'', ``cheapest set of roads
so all cities link'', ``maximise the smallest edge on a path'', ``is this edge
in some / every MST'', ``make the MST unique by changing weights'', ``xor /
special edge weight with an implicit complete graph''}
\recuse{
  \item \textbf{You just need the MST weight or the tree} -- Kruskal
        (\(\secref{c17-kruskal}\)) on sparse, Prim (\(\secref{c17-prim}\)) on
        dense; the cut property is the shared proof.
  \item \textbf{You then ask about paths inside the MST} -- build the
        reconstruction tree (\(\secref{c17-krt}\)) and answer with an LCA weight.
  \item \textbf{The question is about which edges may appear, or uniqueness} --
        the classification (\(\secref{c17-classify}\)) and its equal-weight
        group count.
}
\recnotuse{
  \item \textbf{The answer is a \emph{count} of spanning trees, not a minimum}
        -- that is Ch.30 (\chref{c32}) Matrix-Tree, a determinant.
  \item \textbf{The graph is unweighted and the question is connectivity /
        bridges / articulation} -- that is Ch.29 (\chref{c31}).
  \item \textbf{You need a cheapest \emph{path} from a source, not a cheapest
        spanning tree} -- that is Dijkstra, not Prim (\(\secref{c17-prim}\)).
}
One-line test: \emph{if the input is a weighted graph and the output is
\emph{one} cheapest connected acyclic subgraph (or a question phrased about
that subgraph), you are in this chapter; otherwise the cut is not the tool.}
\end{recognitionbox}

\subsection*{Constraint \(\to\) Technique}

\begin{center}
\renewcommand{\arraystretch}{1.2}
\setlength{\tabcolsep}{3pt}
\begin{tabular}{@{}p{0.34\linewidth} p{0.34\linewidth} >{\raggedright\arraybackslash}p{0.20\linewidth}@{}}
\toprule
\textbf{Constraint signature} & \textbf{Move} & \textbf{Section}\\
\midrule
\(n,m\le 2\cdot 10^5\), sparse, need MST weight & Kruskal sort + DSU & \(\secref{c17-kruskal}\)\\
dense / adjacency matrix, \(n\le 5000\) & Prim, \(\bigO{n^2}\) array frontier & \(\secref{c17-prim}\)\\
MST built, then "minimax edge between \(u,v\)" queries & reconstruction tree + LCA & \(\secref{c17-krt}\)\\
"is this edge in an MST?"; count / unique MST & cycle property per weight group & \(\secref{c17-classify}\)\\
complete graph, edge weight \(=a_i \oplus a_j\) & trie divide-and-conquer & \(\probref{c17-summit}\)\\
\bottomrule
\end{tabular}
\end{center}

\subsection*{Notation Used Only Here}

\([n]=\{1,2,\dots,n\}\) is the vertex set.  A graph \(G=(V,E)\) has \(|V|=n\),
\(|E|=m\), and edge weights \(w\).  A \term{spanning tree} is an acyclic
connected subgraph on all \(n\) vertices, hence exactly \(n-1\) edges; the
\term{minimum spanning tree} minimises \(\sum_{e\in T} w(e)\) over all spanning
trees.  A \term{cut} is a partition \((S, V\setminus S)\); an edge \emph{crosses}
it with one endpoint on each side.  \(T^{*}\) denotes the MST of the running
graph.  For vertices \(u,v\), \(\mathrm{maxw}(u,v)\) is the largest edge weight on the
unique \(u\)-\(v\) path in \(T^{*}\) -- the bottleneck value the reconstruction
tree answers.  \(\oplus\) denotes bitwise xor (used by the Summit).  An edge is
\term{forced} (in every MST), \term{optional} (in some but not all MSTs), or
\term{forbidden} (in no MST).
